\section*{Abstract}

Deploying Vision-Language Models (VLMs) on autonomous mobile robots is severely constrained by the strict computational, memory, and thermal limits of edge hardware. While fine-tuning or retraining foundation models for adverse operational environments is impractical on embedded platforms, compact edge decoders frequently suffer from spatial grounding deficits, resolution downsampling loss, and temporal redundancy in continuous video streams. This thesis investigates an inference-time visual preprocessing architecture designed to enhance VLM scene comprehension and downstream reasoning without modifying underlying model weights.

We propose and evaluate four task-aware spatial prompting techniques Context-Aware Zoom (CAZ), Query-Aware Bounding Box (QABB), Semantic Hazard Isolation (SHI), and Perception Metric Grounding (PMG) alongside dual-track query-aware video keyframe extraction pipelines (QA-IF and Smart Sampling). To systematically benchmark these strategies across the cognitive hierarchy of mobile robotics, we curated a benchmark corpus spanning seven functional task domains, comprising 2,103 static QA pairs across 731 visual scenes (with a 30\% low-light stress test) and 70 continuous robotic video trajectories annotated with minimal keyframe triggers. Real-time tensor inference and telemetry profiling were benchmarked on an NVIDIA Jetson AGX Thor across three representative model tiers: LLaVA-Phi3 (3.8B), Qwen3-VL (4B), and Gemma-4 (12B). To overcome the blind spots of conventional lexical metrics, deep semantic correctness was evaluated via an automated, three-round heterogeneous multi-agent debate jury (DeepSeek-R1, Qwen2.5, and Llama-3) offloaded to an HPC cluster.

Across 12,250 multimodal benchmark instances, our empirical findings reveal that inference-time adaptation fundamentally redefines edge vision-language perception through two governing trade-offs: attentional focus versus global context in static imagery, and extraction latency versus token processing economy in continuous video. While targeted spatial prompts like bounding overlays and zoom crops improve fine-grained object localization by up to +8.81\%, they impose up to a -19.53\% penalty on open-ended scene understanding by stripping peripheral context---demonstrating that visual prompts must be applied conditionally rather than globally. Under adverse sensing conditions, lightweight photometric regularizers such as Gamma Correction reliably stabilize low-light hazard detection without model retraining. In dynamic video streams, while traditional shot detectors fail catastrophically (94.29\% error rate) under smooth robot camera motion, native container demuxing achieves zero-drop real-time throughput (0.45\,s latency), and query-conditioned frame pruning compresses visual payloads down to 1.74 frames to prevent downstream edge VLM token saturation. Ultimately, our analysis uncovers a structural decoupling where standard NLP metrics reward linguistic fluency while overlooking physical hallucinations, establishing that task-conditioned visual preprocessing paired with grounded multi-agent factual auditing provides a practical, robust foundation for embodied edge autonomy.